% =============================================================================
% REMOVED FROM body.tex BY THE 2026-09-12 JOCS REBUILD. NOT \input BY EITHER
% BUILD. DO NOT COMPILE.
%
% WHY THIS FILE EXISTS, AND WHY THE MATERIAL MUST NOT BE PUBLISHED HERE.
% This is the spine of the second paper -- the metadata null across five
% benchmarks plus the NullBench harness -- aimed at the NeurIPS Evaluations &
% Datasets track, next intake ~May 2027 (docs/paper/review/journal_shortlist.md
% Sec. 9, docs/paper/review/null_travels.md, docs/paper/review/nullbench_release.md).
% NeurIPS does not accept work already published in a journal. If any of the
% material below ships inside the JOCS manuscript it is journal-published and
% ED 2027 is closed to it. The split is therefore a publication decision, not a
% length decision, and it is taken deliberately here.
%
% Provenance: every block below is copied verbatim out of docs/paper/body.tex as
% it stood at commit 47dc59d ("Amendment 2: pre-register the least-squares lower
% bound before running it"), with the original line ranges recorded above each
% block.
%
% Labels DEFINED below, which no longer exist in body.tex:
%   sec:splits  sec:covariate  tab:covariate
%
% Labels REFERENCED below that DO still exist in body.tex:
%   sec:interp  sec:limitations  tab:airfrans-sota
% The blocks also reference sec:v2, which the 2026-09-12 rebuild renamed to
% sec:forces; re-point those on reassembly.
%
% One item is deliberately assigned rather than left to land in both papers.
% journal_shortlist.md Sec. 8 flagged the `reynolds`-split finding as genuinely
% ambiguous: it is AirfRANS-specific (so it belongs with the measurement paper)
% but it is also a benchmark-protocol defect (so it strengthens NullBench).
% DECISION: the `reynolds` split goes to PAPER 2, with the null. It is a claim
% about what a published split tests, which is NullBench's subject, and the JOCS
% paper is about measure dependence and representation limits, full stop. The
% BLOCK 3 and BLOCK 5 text below is therefore reserved for paper 2 and is not to
% be reinstated here.
% =============================================================================


% -----------------------------------------------------------------------------
% BLOCK 1 -- Introduction, "The benchmark's force metrics do not measure that
% difference." Original body.tex lines 93-113.
% -----------------------------------------------------------------------------

\paragraph{The benchmark's force metrics do not measure that difference.} AirfRANS
declares the Spearman rank correlation of the force coefficients its primary metric.
Ordinary least squares on the case name, fit on the official $800$ train cases and
scored on the $200$ test cases, ranks the official lift at $\mathbf{0.9821}$
$[0.9737,0.9866]$ and the official drag at $\mathbf{0.9318}$ $[0.8981,0.9531]$. Under the
pre-registered rule---the published mean below the null's point estimate and outside the
null's $95\%$ interval---\textbf{all four AirfRANS-paper lift entries fall below that null}
($0.913$, $0.938$, $0.965$, $0.967$; Transolver's $0.9978$ clears). Measured instead
against each model's \emph{own} published seed spread the separation is decisive only
twice: $3.8$ and $1.9$ seed standard deviations for MLP and PointNet, but $1.6$ and $0.8$
for GraphSAGE and Graph U-Net. We report both readings and rest nothing on the last two.
\textbf{Drag is where the result is not in doubt}: every
AirfRANS-paper entry falls below the null, three of the four published intervals spanning zero
($-0.117$, $-0.022$, $-0.303$, $-0.138$, at published seed standard deviations of
$0.256$, $0.097$, $0.124$, $0.258$). Transolver reports no drag column at all, stating in
its Appendix B.1 that the deep models fail at drag. The reading was fixed in the script's
docstring before any number existed, the in-sample optimism of such a fit is
$0.0008$--$0.0035$, and the two-parameter version of the null is \emph{not}
sufficient---$(U,\alpha)$ alone gives $0.9344$ on lift, which one published entry
straddles. The result needs the digits, and the digits are fair, because they are the
geometry the surrogate is given as an input channel (\autoref{sec:covariate}).


% -----------------------------------------------------------------------------
% BLOCK 2 -- Introduction, "One of the extrapolation splits does not extrapolate."
% Original body.tex lines 115-122.
% -----------------------------------------------------------------------------

\paragraph{One of the extrapolation splits does not extrapolate.} On the
\texttt{reynolds} split, $0$ of $200$ test cases have an inlet velocity inside the
training range---outside on both sides---yet a \emph{nondimensional} parameter
interpolator degrades by $3\%$ on $\texttt{mse\_u}$ and not at all on $\texttt{mse\_p}$,
while the same interpolator on dimensional fields degrades $5.2\times$ and $10.1\times$.
The split's entire difficulty is absorbed by dimensional analysis. The \texttt{aoa}
split, by contrast, bites exactly where it should: $C_l$ relative error rises from
$1.18\%$ to $23.1\%$ (\autoref{sec:splits}).


% -----------------------------------------------------------------------------
% BLOCK 3 -- Introduction, contributions 2 and 3 of six.
% Original body.tex lines 177-186.
% -----------------------------------------------------------------------------

\item \textbf{The benchmark's primary metric is a covariate property of the case index
(\autoref{sec:covariate}).} In-sample optimism for such a fit is $0.0008$--$0.0035$, and
the $(U,\alpha)$ null alone ($0.9344$) is \emph{not} sufficient---the digits are, and they
are fair, being the geometry the surrogate receives as input. We make no leakage
accusation: the claim is that a force-rank number is uninterpretable without this
baseline.

\item \textbf{The \texttt{reynolds} split does not test Reynolds generalisation
(\autoref{sec:splits}).} The difficulty is dimensional scaling, not physics; the
\texttt{aoa} split does bite, and on lift specifically.


% -----------------------------------------------------------------------------
% BLOCK 4 -- sec:splits, in full. Original body.tex lines 1010-1040.
% -----------------------------------------------------------------------------

\subsection{The \texttt{reynolds} split does not test Reynolds generalisation}
\label{sec:splits}

The same estimator, re-selected by cross-validation inside each split's own train set,
separates the benchmark's two extrapolation tasks cleanly. The script proves
disjointness rather than asserting it: on \texttt{reynolds}, \textbf{$0$ of $200$} test
cases have an inlet velocity inside the training range---$U\in[46.8,78.0]$ train
against $[31.5,93.4]$ test, outside on both sides---and on \texttt{aoa},
$0$ of $196$ have an in-range incidence.

Yet the nondimensional interpolator barely notices \texttt{reynolds}: $\texttt{mse\_u}$
degrades by $3\%$ ($0.782\to0.808$) and $\texttt{mse\_p}$ by $-1\%$ ($75.0\to74.1$),
and it is \emph{better} than in-distribution on $\texttt{mse\_v}$ and on surface
pressure. The control that explains why is in the same run: the \texttt{raw}
cell-wise interpolator, identical except that it does not nondimensionalise, degrades
$5.2\times$ on $\texttt{mse\_u}$ ($0.808\to4.194$) and $10.1\times$ on
$\texttt{mse\_p}$ ($74.1\to751.1$) on that split. \textbf{The split's entire difficulty
is absorbed by dimensional analysis, not by learning.} Physically this is unsurprising
---$Re$ spans $2.0$--$6.0\times10^6$, all fully turbulent---but the consequence for
evaluation is sharp: a ``Reynolds OOD'' result on this split is not evidence of learned
extrapolation unless the method is shown not to be exploiting $u/U$ and $p/U^2$.

The \texttt{aoa} split does bite, and only where it should. Lift is the quantity that
depends on incidence, and the $C_l$ relative error rises from $1.18\%$ to
$\mathbf{23.1\%}$; $\texttt{mse\_u}$ degrades $2.1\times$ and $\texttt{mse\_p}$
$1.7\times$, while drag is barely touched ($2.39\%\to2.71\%$). That is the expected
signature of extrapolating past the stall-approach end of the $\alpha$ range. One
caveat is load-bearing: no matched \emph{surrogate} row exists on either split in our
results, so these numbers report the interpolator's own in-distribution-to-OOD
degradation and must not be read as ``interpolation generalises better than a
surrogate''.


% -----------------------------------------------------------------------------
% BLOCK 5 -- sec:covariate, in full, including tab:covariate.
% Original body.tex lines 1042-1136.
% -----------------------------------------------------------------------------

\subsection{The benchmark's force metrics do not measure that difference}
\label{sec:covariate}

\autoref{sec:interp} locates the surrogate's advantage in the first cell off the wall.
The benchmark's declared primary metric cannot see it. \citet{bonnet2022airfrans} name
the Spearman rank correlation of the force coefficients as ``the main metrics to
maximize the recovery of the best airfoils''---and that rank is a covariate property of
the case index.

We fit ordinary least squares on the case name---inlet velocity, incidence, incidence
squared, and the four NACA digits---on the official $800$ train cases, and score
Spearman against the official labels on the $200$ test cases: the exact protocol the
published entries were produced under. The reading was fixed in the script's docstring
before any number existed (\textsc{below the null} iff the published mean is below the
null's point estimate \emph{and} the null's $95\%$ bootstrap interval excludes it).
The null ranks official lift at $\mathbf{0.9821}$ $[0.9737,0.9866]$ and official drag
at $\mathbf{0.9318}$ $[0.8981,0.9531]$ (\autoref{tab:covariate};
\texttt{scripts/covariate\_null\_trainfit.py},
\texttt{results/review/covariate\_null\_trainfit.json}).

\begin{table}[t]
\centering
\small
\caption{The case-name null against the published AirfRANS \texttt{full} field.
Nulls are OLS fit on the official $800$ train cases and scored on the $200$ test
cases; intervals are $95\%$ percentile bootstrap over $10^4$ case-level resamples.
Published rows are from \citet{bonnet2022airfrans} Table 3 and
\citet{wu2024transolver}, quoted with their own seed standard deviations where reported;
Transolver reports no drag column.}
\label{tab:covariate}
\begin{tabular}{l r r}
\toprule
 & $\rho_{C_l}$ vs official & $\rho_{C_d}$ vs official \\
\midrule
null: $\alpha$ & $0.9335$ $[0.9114,0.9483]$ & $0.8001$ $[0.7249,0.8563]$ \\
null: $(U,\alpha)$ & $0.9344$ $[0.9123,0.9492]$ & $0.8276$ $[0.7601,0.8768]$ \\
null: $(U,\alpha,\alpha^2)$ & $0.9344$ $[0.9123,0.9492]$ & $0.8752$ $[0.8278,0.9083]$ \\
\textbf{null: $(U,\alpha,\alpha^2,\text{NACA})$} & $\mathbf{0.9821}$ $[0.9737,0.9866]$ & $\mathbf{0.9318}$ $[0.8981,0.9531]$ \\
\midrule
MLP & $0.913{\pm}0.018$ & $-0.117{\pm}0.256$ \\
GraphSAGE & $0.965{\pm}0.011$ & $-0.303{\pm}0.124$ \\
PointNet & $0.938{\pm}0.023$ & $-0.022{\pm}0.097$ \\
Graph U-Net & $0.967{\pm}0.019$ & $-0.138{\pm}0.258$ \\
Transolver & $0.9978$ & not reported \\
\bottomrule
\end{tabular}
\end{table}

\textbf{Every baseline in the AirfRANS paper sits below the full null on both
coefficients}, on the metric that paper declares primary---four of the five
published entries once Transolver is included, which clears on lift ($0.9978$) and
reports no drag column at all, stating in its Appendix B.1 that ``all the deep models fail
in drag coefficient estimation''. The drag column is worse than the gap alone suggests:
\textbf{three of the four published drag intervals span zero} at one published seed
standard deviation ($-0.117{\pm}0.256$, $-0.022{\pm}0.097$, $-0.138{\pm}0.258$; only
GraphSAGE's $-0.303{\pm}0.124$ excludes it, and on the wrong side). Four qualifications
belong with this, and we state them before a reviewer does.

\emph{On lift the margin is decisive for two entries, not four.} The rule above compares a
published mean against the null's point estimate and its interval, and all four
AirfRANS-paper lift entries fall below it on that rule. A reader who compares instead
against each model's own published spread will not see four clean results: the gap is
$3.8$ published seed standard deviations for MLP ($0.913{\pm}0.018$) and $1.9$ for
PointNet ($0.938{\pm}0.023$), but only $1.6$ for GraphSAGE ($0.965{\pm}0.011$) and $0.8$
for Graph U-Net ($0.967{\pm}0.019$). The two intervals are not commensurable in any case:
ours is a case-level bootstrap over the $200$ test cases, theirs a spread over five
training seeds, and neither carries the other's variance source. The summary we stand
behind is therefore that the null's point estimate exceeds all four lift entries and
separates decisively from two of them. None of this qualification applies to drag, where
the null clears every entry by a wide margin and three of the four intervals span zero.

\emph{The digits are fair.} The NACA digits are the geometry the surrogate is handed as
an input channel, so a null that uses them is not exploiting information the models
lack; it is asking whether the models extract anything from that geometry beyond what a
seven-coefficient regression already extracts. Restricting the null to
$(U,\alpha)$---the reading an earlier version of this control used---gives $0.9344$ on
lift, which PointNet straddles and MLP sits just below at $1.2$ seed standard
deviations. \textbf{The two-parameter null is not sufficient to support this claim; the
digits are what make it decisive.}

\emph{The fit is honest.} The numbers above avoid in-sample optimism by construction,
being fit on train and scored on test. A companion run that fits and scores entirely
within the $200$ test cases bounds the optimism directly: the gap between the in-sample
and the 10-fold cross-validated Spearman is $0.0017$ on lift and $0.0020$ on drag at
the full feature set, and $0.0008$--$0.0035$ across all four feature sets
(\texttt{scripts/covariate\_null.py}, \texttt{results/review/covariate\_null.json}).
Covariate regressions of this form carry negligible optimism at this sample size.

\emph{This is a statement about the metric, not an accusation.} We do not claim label
leakage. A drag \emph{rank} metric cannot distinguish a better integrator from a head
that has learned the covariate dependence, and conditioning a surrogate on its inlet
state is necessary rather than suspect. The defensible claim is the measurement gap: no
paper in this line reports the null, so its force-rank numbers are uninterpretable as
evidence about integration quality. The claim is also narrower than ``nobody publishes
parameter-space baselines'', which is false as stated (\autoref{sec:interp}).


% -----------------------------------------------------------------------------
% BLOCK 6 -- the covariate half of the drag-integrator paragraph. Original
% body.tex lines 1200-1211, inside "What these force correlations measure".
% The rebuilt manuscript KEEPS that paragraph but stops at the integrator's own
% limits; the three sentences below, which make the covariate argument, moved
% here with the null.
% -----------------------------------------------------------------------------

Nor is it limited only by
the measurement: a three-parameter OLS on $(U,\alpha,\alpha^2)$ that never touches a flow
field ranks the official drag at $\rho=0.874$---above every integrator arm we tried on
ground-truth fields ($0.611$--$0.839$) and on predicted fields
($0.828$--$0.845$); $\alpha$ alone gives $0.800$
(\texttt{results/control/drag\_covariate\_control.json}). The parameter interpolator makes
the same point from the other direction: scored through this integrator it ranks the
official drag at $\rho=0.8389$, \emph{indistinguishable from the same integrator applied
to the exact ground-truth field} ($0.8394$), and matching it on lift to four significant
figures ($0.8758$ vs $0.8757$). \textbf{A drag rank correlation reported on this benchmark
is uninterpretable without its covariate baseline}, and the conventional near-field number
sits below that baseline.


% -----------------------------------------------------------------------------
% BLOCK 7 -- the covariate and split limitations bullets.
% Original body.tex lines 1854-1859 and 1864-1869.
% -----------------------------------------------------------------------------

\item \textbf{No matched surrogate row on the extrapolation splits.} The degradations in
\autoref{sec:splits} are the interpolator's own in-distribution-to-OOD changes, so they
must not be read as ``interpolation generalises better than a surrogate''. Two scope notes:
the \texttt{reynolds} test split is the first $200$ of $496$ cases (the $r128$ cache
limit), an unbiased subsample whose $U$/$\alpha$ ranges match the full set but a subsample
nonetheless; and \texttt{aoa} has $196$ test cases, not $200$.

\item \textbf{The covariate null is one benchmark, one metric family.} We compute it for
AirfRANS lift and drag only. Extending it to DrivAerNet++, DrivAerML and AhmedML is
straightforward---all three publish their design parameters---and is not done here. We also
do not refit the null inside the \texttt{scarce}, \texttt{reynolds} and \texttt{aoa}
splits, which would be required before comparing published per-split force numbers against
it.


% -----------------------------------------------------------------------------
% BLOCK 8 -- the covariate sentences of the Conclusion.
% Original body.tex lines 1971-1978.
% -----------------------------------------------------------------------------

Meanwhile a regression on the case name outranks every published drag
entry, three of which span zero, and outranks all four published lift entries as
well---though on lift the separation is $3.8$ and $1.9$ of each model's own published seed
standard deviations for two of them and only $1.6$ and $0.8$ for the other two. We do not
read this as evidence that those results are wrong; we read it as evidence that both a
field metric and a force-rank metric on this benchmark are uninterpretable without their
non-learned reference point---and, for the field metric, without the weighting it is
computed under.

% and, from the closing recommendation paragraph (original line 2000-2001):

Anyone reporting a force-rank metric should report
it against a regression on the case name, which takes seconds.
